% Solutions to the exercises of Chapter 6.
\begin{solutions}
  \item With $x_k = \sum_{i<k}w_i$ and $x_j = \sum_{l<j}w_l$,
        \[
          C(k, j) = \E[x_kx_j] = \sum_{i<k}\sum_{l<j}\E[w_iw_l] = \sum_{i<k}\sum_{l<j}\sigma^2\delta_{il} = \sigma^2\sum_{i<\min(k, j)}1 = \min(k, j)\,\sigma^2 .
        \]
        Only the shared terms of the two sums are correlated; the later sum contains the whole earlier one, which is why a random walk
        ``remembers'' its past.

  \item The stationary variance is $\sigma^2/(1 - a^2) = 1/0.0199 = 50.3$, a standard deviation of $7.09$: an input of unit size per step builds up a
        seven times larger wandering error. The correlation $a^{|j|}$ drops below $1/e$ when $|j| > 1/(-\ln a) = 99.5$, i.e.\ after 100 steps,
        consistent with the correlation time $1/(1 - a) = 100$ steps.

  \item By \eqref{eq:discrete-noise} a single sample at $\Dt = 1$~ms has standard deviation $\sigma/\sqrt\Dt = 2\times10^{-3}\sqrt{1000} = 0.063~\mathrm{m/s^2}$.
        The average of 1000 consecutive samples (one second) has $0.063/\sqrt{1000} = 2\times10^{-3}~\mathrm{m/s^2}$, which is $\sigma/\sqrt T$ with $T = 1$~s:
        averaging over a time $T$ always gives $\sigma/\sqrt T$, however the samples are spaced. Read at 100~Hz with its anti-aliasing filter set
        accordingly, each sample averages 10~ms and has $\sigma\sqrt{100} = 0.02~\mathrm{m/s^2}$. (If the 100~Hz stream were obtained by merely keeping
        every tenth 1~kHz sample, each sample would keep the full $0.063~\mathrm{m/s^2}$, and high-frequency noise would alias into the band.)

  \item The alias frequency is $|f - kf_s|$ for the integer $k$ that brings it into $[0, f_s/2] = [0, 50]$~Hz: $|105 - 100| = 5$~Hz. A propeller
        vibration just above the sampling rate appears as a slow 5~Hz oscillation, indistinguishable from real motion.

  \item Simulate $R = 1000$ runs of white noise with per-sample standard deviation $\sigma/\sqrt\Dt$ at $\Dt = 5$~ms, integrate by cumulative sums
        (\code{np.cumsum(w, axis=1) * dt} once, twice and three times), and compute the sample variance of each integral at $t = 10$~s. One run of
        this experiment with $\sigma = 1$ gives
        \begin{center}\small
        \begin{tabular}{@{}l l l l@{}}
        \toprule
        $m$ & Monte Carlo variance & theory $\sigma^2t^{2m-1}/((2m-1)((m-1)!)^2)$ & ratio \\
        \midrule
        1 & 9.78 & 10 & 0.978 \\
        2 & 327.2 & 333.3 & 0.982 \\
        3 & 4964 & 5000 & 0.993 \\
        \bottomrule
        \end{tabular}
        \end{center}
        With 1000 runs the variance ratio has a 95\% band of $\pm1.96\sqrt{2/999} = \pm8.8\%$ (\cref{sec:chi2}), so all three agree with
        \cref{prop:integrated-noise}. The discretization error of the rectangle rule is $O(\Dt/t) \approx 0.05\%$, far below the sampling error.

  \item With $x = (p, v)$, $F = \begin{bmatrix}0 & 1\\0 & 0\end{bmatrix}$, $G = \begin{bmatrix}0\\1\end{bmatrix}$ and $Q_c = \sigma^2$, the three distinct entries of
        $\dot P = FP + PF\T + GQ_cG\T$ are
        \[
          \dot P_{pp} = 2P_{pv}, \qquad \dot P_{pv} = P_{vv}, \qquad \dot P_{vv} = \sigma^2 .
        \]
        Integrating from $P(0) = 0$ in cascade: $P_{vv} = \sigma^2t$, $P_{pv} = \sigma^2t^2/2$, $P_{pp} = \sigma^2t^3/3$, in agreement with $m = 2$ of
        \cref{prop:integrated-noise}. The correlation coefficient $P_{pv}/\sqrt{P_{pp}P_{vv}} = \sqrt3/2$ is constant: position and velocity errors are
        strongly correlated, because the position error is the integral of the velocity error.

  \item For $\dot x = -x/\tau + w$, the Lyapunov equation is $\dot P = -2P/\tau + \sigma^2$, whose solution is
        \[
          P(t) = e^{-2t/\tau}P(0) + \frac{\sigma^2\tau}{2}\big(1 - e^{-2t/\tau}\big),
        \]
        approaching $P_\infty = \sigma^2\tau/2$. For the discretization, the scalar version of \cref{thm:lti-solution} gives
        $x(t + \Dt) = e^{-\Dt/\tau}x(t) + \int_0^{\Dt}e^{-(\Dt - s)/\tau}w(t + s)\,\dd s$. The noise integral is independent of $x(t)$, with variance
        $\sigma^2\int_0^\Dt e^{-2(\Dt - s)/\tau}\,\dd s = \frac{\sigma^2\tau}{2}\big(1 - e^{-2\Dt/\tau}\big) = P_\infty\big(1 - e^{-2\Dt/\tau}\big)$. Hence
        $x_{k+1} = e^{-\Dt/\tau}x_k + w_k$ with $\Var(w_k) = P_\infty(1 - e^{-2\Dt/\tau})$, exact for any $\Dt$. Two checks: if $x_k$ has variance $P_\infty$, so does
        $x_{k+1}$ (stationarity is preserved); and for $\Dt \ll \tau$, $\Var(w_k) \approx \sigma^2\Dt$, the random-walk increment.

  \item Simulate the gyro as white noise of density $\sigma_g$ plus a random walk of density $\sigma_{b_g}$ (\cref{sec:ct-noise} gives both per-sample
        conversions), integrate the rate to an angle $\theta$, and evaluate the overlapping Allan variance
        $\sigma_A^2(\tau) = \langle(\theta_{i+2m} - 2\theta_{i+m} + \theta_i)^2\rangle/(2\tau^2)$, $\tau = m\Dt$, on a log-spaced grid of $\tau$ up to a quarter of the
        record. For a 2-hour record the curve has the two predicted slopes, with its minimum near $\tau = 16$~s, close to the crossover
        $\tau = \sqrt3\,\sigma_g/\sigma_{b_g} = 15$~s where the two terms are equal. Fitting $\sigma_A^2 = \sigma_g^2/\tau + \sigma_{b_g}^2\tau/3$ (in log space, weighting each
        $\tau$ by $\sqrt{T/\tau}$, since longer clusters are fewer and noisier) recovers $\sigma_g$ to within about 1\% and $\sigma_{b_g}$ to within 5--10\%.
        Repeating the experiment many times gives the scatter of the estimate of $\sigma_{b_g}$ as a function of record length:
        \begin{center}\small
        \begin{tabular}{@{}l c c c c c@{}}
        \toprule
        record length & 0.5~h & 1~h & 2~h & 4~h & 8~h \\
        \midrule
        relative standard deviation of $\hat\sigma_{b_g}$ & 13\% & 9\% & 6\% & 4.5\% & 3\% \\
        \bottomrule
        \end{tabular}
        \end{center}
        About one hour is needed for 10\% (one sigma). The reason is simple: the random walk dominates only for $\tau \gtrsim 50$~s, and an Allan
        variance computed from $K$ independent clusters has a relative error of about $\sqrt{2/K}$ in variance, half that in deviation; reaching
        10\% in deviation from a single $\tau$ requires $K \approx 50$ clusters of a hundred seconds or more, and combining several $\tau$ lowers this
        somewhat. Real sensors add bias instability and temperature drift, which is why
        calibration recordings usually last several hours.
\end{solutions}
